\SourceLeaf{522}Damit entwickelt sich $\Sigma[-\mathfrak{Q}]\times\Sigma[\mathfrak{Q}]$ zu
\[
\begin{aligned}
&-[0]+[1]+[4]+[9]\dotsfive+[(n-1)^2]\\
&+[0]+[3]+[8]+[15]\dotsfive+[nn-1]\\
&+[0]+[5]+[12]+[21]\dotsfive+[nn+2n-3]\\
&+[0]+[7]+[16]+[27]\dotsfive+[nn+4n-5]\\
&+\text{usw.}\\
&+[0]+[2n-1]+[4n]+[6n+3]\dotsfive+[3nn-6n+3]
\end{aligned}
\]
Addiert man diese Teile \textit{spaltenweise}, so erhält man
\[
\begin{aligned}
&n[0]+[1]\times\frac{1-[2n]}{1-[2]}+[4]\frac{1-[4n]}{1-[4]}+[9]\frac{1-[6n]}{1-[6]}\\
&\qquad+\text{usw.}+[(n-1)^2]\times\frac{1-[2nn-2n]}{1-[2n-2]}
\end{aligned}
\]
In diesem Ausdruck verschwinden alle Teile außer dem ersten, wenn $n$ ungerade ist; denn dann werden $1-[2n]$, $1-[4n]$, $1-[6n]$ usw. sämtlich $=0$, dagegen keiner der Nenner $1-[2]$, $1-[4]$, $1-[6]$, $1-[8]$ usw. bis $1-[2n-2]$. Ist $n$ hingegen gerade, so ist auch einer der Nenner $=0$, nämlich $1-[n]$, dem der Term $[\frac14nn]\times\frac{1-[nn]}{1-[n]}$ entspricht; die Summe der Teile, aus denen dieser entstanden ist, wird jedoch $=n[\frac14nn]$. Hier sind erneut zwei Fälle zu unterscheiden. Ist $n$ durch vier teilbar, so gilt $\frac14nn\equiv0\, (\mathrm{mod}.n)$ und daher $[\frac14nn]=1$; ist $n$ dagegen das Doppelte einer ungeraden Zahl, so gilt $\frac14nn\equiv\frac12n\, (\mathrm{mod}.n)$ und daher notwendig $[\frac14nn]=-1$. Schließlich folgt daraus
\par\noindent 1) für einen ungeraden Wert von $n$ ist das gesuchte Produkt $=n$
\par\noindent 2) für einen durch vier teilbaren Wert ist das Produkt $=2n$
\par\noindent 3) für einen Wert, der das Doppelte einer ungeraden Zahl ist, ist es $=0$.\quad \mbox{Q. E. I.}
\par\medskip\Needspace{4\baselineskip}\pdfbookmark[1]{371.}{sectio-octava-art-371}\centerline{\textbf{371.}}\nobreak\smallskip
Es lohnt sich nun, die Beschaffenheit der Summe $\Sigma[\mathfrak{Q}]$ näher zu betrachten.
\par I. Da man die Quadrate $0$, $1$, $4$, $9$, $16$ usw. durch ihre kleinsten Reste modulo $n$ ersetzen darf, gilt offenbar: Bezeichnet $M$ ohne nähere Festlegung die quadratischen Reste der Zahl $n$ von $0$ bis $n-1$ und $m$ die Anzahl der Wurzeln der Kongruenz $xx\equiv M(\mathrm{mod}.m)$, so ist $\Sigma[\mathfrak{Q}]=\Sigma m[M]$. Wie die Zahl $m$ zu bestimmen ist, haben wir in den Artikeln $104$, $105$ gezeigt.
\par II. Ist $n$ eine (ungerade) Primzahl, so gilt für $M=0$ die Gleichung $m=1$, für jeden anderen Wert von $M$ hingegen $m=2$. Ist $n$ eine Potenz einer ungeraden Primzahl $=p^\nu$, so gilt $m=2$ für jeden Wert von $M$, der nicht durch $p$ teilbar ist\quad\textemdash
\par\nobreak\noindent\hbox to\linewidth{\textemdash\hfil\textemdash\hfil\textemdash\hfil\textemdash\hfil\textemdash\hfil\textemdash\hfil\textemdash\hfil\textemdash\hfil\textemdash\hfil\textemdash\hfil\textemdash\hfil\textemdash\hfil\textemdash\hfil\textemdash\hfil\textemdash\hfil\textemdash\hfil\textemdash}\par
